%!TEX root = ../Thesis.tex
\chapter{Plan}


\section{Gantt Chart}

\noindent
Figure~\ref{fig:gantt} shows an outline of the project, including the studying, creation, and evaluation of the pipeline in question.

\begin{sidewaysfigure}
    \centering
    \includegraphics[width=\textheight]{gantt/gantt.pdf}
    \caption{Project timeline (Gantt chart)}
    \label{fig:gantt}
\end{sidewaysfigure}

\section{Activities}

\note{This section provides a table of the main phases and activities of the plan with a little description of each.}

\begin{table}[h]
\caption{Project phases and activities}
\centering{
\scriptsize{
\begin{tabular}{|c p{2.0cm}| p{5.0cm}| p{2.0cm} | c|}
\toprule 
\hline
\textbf{\#}	& \textbf{Title} &	\textbf{Description} &	\textbf{Phase} &	\textbf{Period}  \\
\hline
\hline
$15$	& Project Definition   &	Writing up the PDR and planning the project. Also includes attending a start-up supervisor meeting. & 	 Project Definition & W1--W2 \\
$23$	& Writing up Thesis    &	This is the final writing up -- all of the chapters should be ready by now & 	 Thesis & W26	 \\
\hline
\bottomrule
\end{tabular}
}
}
\label{tab:survey}
\end{table}


\section{Milestones}

\note{This section provides a table of the main milestones of the plan with a little description of each. Examples of milestones include: 
(i) Approvals (e.g., of a UX design, data management, or of a study involving patients); 
(ii) Important deadlines (e.g., prototype ready for UX testing, recruitment of all patients, or final analysis of data done); 
(iii) Acquisition of core resources (e.g., getting access to hardware, setting up a lab experiment, or getting access to test data); and 
(iv) Thesis milestones (e.g., the \ac{PDR}, chapters, and final Thesis).}


\begin{table}[h]
\caption{Milestones. A: Approval, D: Deadline, AC: Acquisition, TH: Thesis}
\centering{
\scriptsize{
\begin{tabular}{|c p{2.0cm}| p{5.0cm}| c | c|}
\toprule 
\hline
\textbf{\#}	& \textbf{Title} &	\textbf{Description} &	\textbf{Type} &	\textbf{Week}  \\
\hline
\hline
$M_1$	& Project Definition Report    &	Approval of the PDR. & 	 TH & 	 4 \\
$M_2$	& Milestone 2                  &	This is the second milestone.... & 	 D & 	 5 \\
$M_x$	& Milestone x              &	This the x'th milestone.... & 	 D & 	 N \\
$F$	& Thesis hand-in               &	Hand-in of final Thesis & 	 TH & 	 T \\
\hline
\bottomrule
\end{tabular}
}
}
\label{tab:survey}
\end{table}


\section{Deliverable}

\note{This section provides a table of the main deliverables of the plan with a little description of each.
Examples of deliverables include:
(i) Design products (e.g., UX design, a requirement analysis, a design model, a data model), 
(ii) Hardware and Software (e.g., a UX prototype, a lab setup, an app for download in App Store);
(iii) Data analysis (e.g., machine learning analysis or statistical analysis of user study data); and
(iv) Thesis deliverables (e.g., chapters, figures).}


\begin{table}[h]
\caption{Deliverable. R: Report, SW: Software, ST :Study}
\centering{
\scriptsize{
\begin{tabular}{|c p{2.0cm}| p{5.0cm}| c | c|}
\toprule 
\hline
\textbf{\#}	& \textbf{Title} &	\textbf{Description} &	\textbf{Type} &	\textbf{Week}  \\
\hline
\hline
$D_1$	& Project Definition Report    &	Final version of the PDR. & 	 R & 	 4 \\
$D_2$	& Milestone 2                  &	This is a software milestone.... & 	 SW & 	 8 \\
$D_2$	& Final Thesis                 &	Final proof-read Thesis ready to hand in & 	 R & 	 T \\
\hline
\bottomrule
\end{tabular}
}
}
\label{tab:survey}
\end{table}


\section{Risk Analysis}

\note{A risk analysis identifies the top-5 or so risks that can significantly upset the entire project. Typical risk factors are often those that are not under your own control. Examples include getting access to data from someone else, the ability to run some experiments that will generate data, getting access to hardware/software that is needed, and getting access to people, such as company representatives, users, and patients.}


\begin{table}[h]
\caption{Risks. 1:low risk, 5:high risk of delays}
\centering{
\scriptsize{
\begin{tabular}{|c p{2.0cm}| p{5.0cm}| c | c|}
\toprule 
\hline
\textbf{\#}	& \textbf{Title} &	\textbf{Description} &	\textbf{Risk} &	\textbf{Week}  \\
\hline
\hline
$R_1$	& Risk no. 1    &	This is a high risk. Describe here how to mitigate it. & 	 5 & 	 4 \\
$R_2$	& Risk no. 2    &	This is a low risk. Describe here how to mitigate it. & 	 2 & 	 4 \\
\hline
\bottomrule
\end{tabular}
}
}
\label{tab:survey}
\end{table}

